\chapter{Legal and Ethical Aspects of Genomic Data Use}

\section{Introduction}

The contemporary genomic revolution, characterized by exponential reduction in sequencing costs and unprecedented increase in availability of DNA and proteome data at population scale, presents profound legal, ethical, and regulatory challenges. The contemporary researcher in phylogenomics and bioinformatics operates in a complex context marked by multiple regulatory frameworks governing the collection, storage, sharing, and use of genomic data. These frameworks vary substantially between national jurisdictions, often in tension with scientific norms of data openness and reproducibility.

This chapter examines the legal and ethical aspects central to contemporary genomic research, with particular emphasis on the Brazilian context. We cover the international treaties establishing the sovereign right of States over genetic resources within their borders; the specific Brazilian legislation regulating access to biodiversity, protection of personal data, and ethical use of animals in research; the international regulatory frameworks shaping the sharing of genomic data; and the emerging ethical considerations in genomic research with indigenous and traditional communities.

\section{International Frameworks: Convention on Biological Diversity and Nagoya Protocol}

\subsection{The Convention on Biological Diversity (1992)}

The Convention on Biological Diversity (CBD), adopted at the United Nations Conference on Environment and Development in Rio de Janeiro in 1992 and in force since December 1993, establishes an international legal framework for the conservation of biological diversity, the sustainable use of its components, and the fair and equitable sharing of benefits derived from the use of genetic resources \parencite{cbd1992}. The CBD recognizes the sovereignty of States over genetic resources in their territories and establishes that access to genetic resources is a matter of internal competence of each contracting party.

The key principles of the CBD include: (1) recognition of state sovereignty over genetic resources within national jurisdictions; (2) right of States to legislate on access to genetic resources within their territories; (3) obligation to facilitate access to genetic resources on mutually agreed terms; and (4) obligation to share benefits derived from the use of genetic resources in a fair and equitable manner. These principles form the foundation for all subsequent legislation on access to biodiversity at the global level.

For researchers, the CBD implies that collection of biological specimens or genomic data from organisms in a country requires prior informed consent (PIC) from the competent authorities of that country and results in obligations for benefit-sharing \parencite{cbd1992}.

\subsection{Nagoya Protocol on Access and Benefit Sharing (2014)}

The Nagoya Protocol on Access and Benefit Sharing, adopted in 2010 and in force since October 2014, operationalizes and provides specific mechanisms for implementation of the principles established by the CBD \parencite{nagoyaprotocol2014}. The protocol requires that signatory countries establish legislative, administrative, and policy measures to: (1) ensure access to genetic resources conditioned on prior informed consent; (2) require disclosure of the origin of genetic resources in applications for intellectual property rights; (3) establish requirements for mutually agreed terms documenting the conditions of access; and (4) monitor the use of genetic resources along supply chains.

For biologists and researchers in genomics, the Nagoya Protocol means that genomic data from organisms collected in signatory countries cannot be freely shared or published without attention to the benefit-sharing obligations established through mutually agreed terms. This requirement may impose substantial restrictions on data publication and the open science practice that has historically been the norm in genomics.

\section{Brazilian Legislation: Biodiversity and Data Protection}

\subsection{Biodiversity Law (Lei nº 13.123/2015) and SisGen}

Brazil, holder of approximately 13\% of global biodiversity and signatory to both the CBD and the Nagoya Protocol, implemented the Biodiversity Law (Lei nº 13.123, of May 20, 2015) and subsequently Decree nº 8.772/2016, which regulate access to components of Brazilian biodiversity and the sharing of benefits derived from such access \parencite{brasil2015, brasil2016}.

The Biodiversity Law establishes that any access to a component of Brazilian biodiversity for purposes of scientific research, bioprospecting, conservation, or sustainable use requires prior registration in the Genetic Heritage Management System (SisGen). SisGen is an electronic information system that consolidates data on access to components of Brazilian genetic heritage, derived products, and reproductive material obtained from access, with the purpose of: (1) identifying the origin of genetic materials; (2) tracking the access history; (3) facilitating benefit-sharing; and (4) complying with Brazil's obligations under the Nagoya Protocol.

Researchers in phylogenomics must register with SisGen before accessing, collecting, or using components of Brazilian biodiversity for research purposes. This includes collection of specimens for sequencing, obtaining tissues from organisms of Brazilian origin, and analysis of genomic sequences derived from Brazilian biodiversity. Failure to register with SisGen constitutes an environmental violation and can result in fines of up to R\$ 200,000, in addition to possible revocation of research funding and publication restrictions \parencite{brasil2015}.

A critical aspect of the Biodiversity Law is that it classifies access to biodiversity components into three categories: (1) scientific research without economic purposes; (2) bioprospecting; and (3) sustainable use. Researchers in phylogenomics frequently work in the category of scientific research without economic purposes, but if the research results in a product of economic interest (for example, identification of a gene with biotechnological applications), the classification may change retroactively, triggering additional benefit-sharing obligations.

\subsection{General Data Protection Law (LGPD - Lei nº 13.709/2018)}

The General Data Protection Law (LGPD), promulgated in August 2018 and in full effect since August 2020, is the Brazilian legislation for personal data protection \parencite{brasil2018lgpd}. Although often considered analogous to the General Data Protection Regulation (GDPR) of the European Union, the LGPD is specific Brazilian legislation that establishes rights and responsibilities for the processing of personal data in Brazil.

The LGPD is relevant to genomic research because genomic data, when associated with an identifiable or potentially identifiable individual, are classified as sensitive personal data. The law defines personal data as ``information related to an identified or identifiable natural person'' \parencite{brasil2018lgpd}. Genomic data are often considered a special category of sensitive data because they can reveal information about health, genetic predisposition to disease, or ancestry.

Implications of the LGPD for genomic research include: (1) explicit consent required for collection and processing of personal genomic data; (2) right of individuals to access, correct, and delete their genomic data; (3) obligation to implement technical and administrative security measures for data protection; (4) obligation to report data breaches within a specified timeframe; and (5) right of individuals to object to processing of their genomic data for specific purposes \parencite{brasil2018lgpd}.

For researchers, the LGPD means that any research project involving genomic data from Brazilian humans requires: written informed consent from the donor or their legal representatives; transparent privacy policy documenting how data will be collected, processed, stored, and shared; appropriate security measures; and a data retention plan documenting how long data will be maintained.

\subsection{Brazilian Legislation on Research with Animals}

For phylogenomic research on metazoans, Brazilian legislation on ethical use of animals in research is also relevant. The Arouca Law (Lei nº 11.794, of October 8, 2008) regulates the scientific use of animals in teaching and research in Brazil \parencite{brasil2008arouca}. The law establishes that: (1) every institution that conducts teaching or research activities with animals must establish an Animal Use Ethics Committee (CEUA); (2) every research proposal involving animals must be submitted to the CEUA before beginning; (3) the CEUA must evaluate whether the research is scientifically appropriate, ethically justified, and uses techniques that minimize animal suffering.

Researchers in phylogenomics who collect live animal specimens, who conduct experiments with animals, or who use animal tissues obtained through euthanasia for collection purposes must submit their protocols to the CEUA of their institution. The Arouca Law establishes principles of: (1) use of alternatives to animals when available; (2) reduction in the number of animals used (the 3Rs principle: Reduction, Replacement, Refinement); (3) refinement of techniques to minimize pain and suffering; and (4) humane euthanasia when necessary \parencite{brasil2008arouca}.

\section{International Frameworks for Genomic Data Sharing}

\subsection{International Sequence Infrastructure: GenBank, ENA, and DDBJ}

The international infrastructure for genomic sequence data rests on three central databases linked globally that share data daily: the National Center for Biotechnology Information (NCBI) GenBank maintained by the National Institutes of Health (NIH) of the United States, the European Nucleotide Archive (ENA) of the European Bioinformatics Institute, and the DNA Data Bank of Japan (DDBJ) \parencite{genbank2023, ena2023, ddbj2023}. Collectively, this infrastructure is referred to as the International Nucleotide Sequence Database Collaboration (INSDC).

The INSDC databases operate under the principle of immediate public access: sequences deposited are publicly accessible without restriction within 24 hours of submission, unless the submitter requests temporary retention of up to 4 months (called a hold period). This system was designed to facilitate rapid scientific discovery and ensure reproducibility of published analyses, allowing researchers to reanalyze underlying data for published studies.

However, for genomic data from humans that may be potentially identifying, the INSDC databases have implemented access restrictions. GenBank maintains a separate database for human genomic data that requires consent from the proposer for any public release, allowing indefinite retention of sensitive data while keeping the sequence available to researchers approved through a controlled access portal.

For researchers in phylogenomics working with metazoans, the implication is that sequences from non-human organisms are routinely deposited in INSDC with immediate public access. However, if the research involves genomic data from humans or specific human populations, the researcher must be aware of the Brazilian LGPD, the European GDPR, and other data protection laws that may impose restrictions on public data sharing.

\subsection{Open Data Mandate Policies of Funding Agencies}

International research funding agencies have increasingly adopted mandatory data sharing policies that require publicly funded researchers to make their data available in public repositories within a specified timeframe (often 6 months to 2 years after publication of results). The National Institutes of Health (NIH) of the United States, the European Research Council (ERC), and the Engineering and Physical Sciences Research Council (EPSRC) of the United Kingdom have all implemented open data policies.

In Brazil, funding agencies such as the São Paulo Research Foundation (FAPESP) and the National Council for Scientific and Technological Development (CNPq) have adopted similar policies, often in conformity with UNESCO recommendations for open data for science \parencite{unesco2021}. Funding for research by FAPESP or CNPq now typically requires that raw sequencing data be deposited in recognized public repositories and that analyses be reproducible through sharing of code and methods.

These policies have significant implications for phylogenomic research: genomic data must be planned to be shareable from the initial study design. For sensitive data (humans, data from indigenous populations), controlled access policies may be appropriate, but must be negotiated with funding agencies during proposal development.

\section{Ethical Considerations in Genomic Research with Indigenous and Traditional Communities}

\subsection{Rights of Indigenous Communities and Biopiracy}

Brazil is a signatory to the United Nations Declaration on the Rights of Indigenous Peoples (UNDRIP), which establishes collective rights of indigenous peoples over their traditional knowledge, cultural and genetic resources \parencite{undrip2007}. UNDRIP establishes that indigenous peoples have the right to: (1) consent or not consent to the collection and use of their biodiversity and associated knowledge; (2) control, protect, and preserve their traditional knowledge; and (3) maintain and strengthen their distinctive relationship with lands and resources.

Historically, genetic resources associated with traditional knowledge of indigenous communities were collected and used in research and commercial development without consent or benefit to indigenous communities, a process often termed ``biopiracy.'' Notable cases include collection of medicinal plants and organisms from the Brazilian Amazon used in pharmaceutical development without participation or financial benefit to indigenous communities.

Brazilian Biodiversity Law establishes that access to traditional knowledge associated with a biodiversity component requires prior informed consent from the indigenous community or traditional population and results in benefit-sharing obligations. SisGen requires explicit documentation of consent when research accesses knowledge associated with biodiversity components.

For researchers in phylogenomics, this means that if the research involves collection of organisms on indigenous lands, or uses genomic data from organisms whose traditional knowledge was used in the research or development, the researcher must: (1) obtain prior informed consent from the indigenous community; (2) document consent and benefit-sharing in SisGen; (3) ensure that communities participate in decisions about data publication and sharing; and (4) implement benefit-sharing mechanisms, which may include financial participation or academic collaboration.

\subsection{Principles of Respectful Research with Indigenous Communities}

Beyond legal requirements, there is an ethical imperative to conduct research with indigenous communities in a respectful and participatory manner. Organizations such as the International Society of Ethnobiology have established ethical principles for research in ethnobiology, ethnobotany, and ethnogenomics \parencite{ise2006}.

These principles include: (1) reciprocity --- research should benefit collaborating communities, not just researchers; (2) transparency --- communities must fully understand research objectives and possible uses; (3) sovereignty --- indigenous communities must maintain control over their genomic data and traditional knowledge; (4) respect --- researchers must respect indigenous cultural protocols and intellectual rights; and (5) benefit-sharing --- researchers must share financial and non-financial benefits with communities.

In practice, this means that research involving genomics of organisms associated with indigenous traditional knowledge must include: (1) prior dialogue with communities to explain research objectives; (2) documented consent obtained freely without coercion; (3) participation of community members in research design when possible; (4) mechanisms for sharing results in the community's language and in accessible formats; and (5) pre-agreed benefit-sharing if research results in commercial developments.

\section{European GDPR and International Genomic Data Protection}

\subsection{General Data Protection Regulation (GDPR) and Genomic Sequences}

The General Data Protection Regulation (GDPR), adopted by the European Union in 2016 and in force since May 2018, is data protection legislation even more stringent than the Brazilian LGPD. The GDPR is particularly relevant to genomic research because: (1) it classifies genomic data as a special category of personal data requiring explicit consent; (2) it establishes individuals' right of access to their genomic data; (3) it establishes the right to delete data (``right to be forgotten''); and (4) it requires that any processing of genomic data be proportionate and justified \parencite{gdpr2018}.

For researchers working with genomic data from European individuals, the GDPR imposes substantial obligations. Genomic data cannot be transferred outside the European Union without guarantees of data protection equivalent to GDPR. Researchers need a legal basis for processing genomic data (typically informed consent) and must document this legal basis in a Data Protection Impact Assessment (DPIA).

\subsection{Implications for International Genomics Research}

Researchers in phylogenomics working in an international context must be aware that different jurisdictions involved in research may have different requirements. A collaborative project involving collection of genomic data in Brazil, processing and analysis in Europe, and publication in an American journal may be simultaneously subject to: Brazilian Biodiversity Law and SisGen; Brazilian LGPD if data are from Brazilian humans; European GDPR if analysis occurs in a European institution; and United States data export regulations if data are transferred to American institutions.

Compliance with all these frameworks requires deliberate planning from the initial study design. Recommendations include: (1) consult regulatory compliance personnel and legal advisors during study design; (2) document the origin, consent, and use restrictions of all data from the point of collection; (3) implement appropriate technical data security measures; (4) establish international data transfer agreements if data will be shared between jurisdictions; and (5) maintain a record of all regulatory authorities involved and their specific requirements.

\section{Open Access and Open Science in Genomics}

\subsection{Movement for Open Data and Open Science}

There is a growing global movement toward what is often called ``open science'' or ``open data'' --- the practice of making raw scientific data, computational code, and analyses publicly available in a manner that any researcher can access, review, and reuse this data and code without restriction. UNESCO published a Recommendation on Open Science in 2021, endorsed by Brazil, that establishes principles of openness in scientific data \parencite{unesco2021}.

For genomics, the open data movement has been particularly influential. The Human 1000 Genomes Project, completed in 2015, set a powerful precedent by making publicly accessible, without restriction, the genomic data of 2,500 individuals from diverse global populations. This open access to data allowed thousands of researchers to use the data for their own research, substantially accelerating progress in population genomics and precision medicine \parencite{1000genomes2015}.

However, the open data model for genomics is in tension with personal data protection legislation such as GDPR and LGPD, which impose restrictions on public sharing of genomic data from identifiable humans. Resolution of this tension continues to evolve, with ongoing discussions about how to implement ``responsible access'' --- an intermediate model in which genomic data are publicly accessible for scientific research through authorization processes, but with restrictions on commercial sharing or re-identification \parencite{ga4gh2018}.

\subsection{Controlled Access Repositories and Data Governance Models}

To resolve the tension between open data and personal data protection, controlled access repository infrastructure and data governance models have been developed. The Global Alliance for Genomics and Health (GA4GH) has established international standards for federating controlled access genomic databases \parencite{ga4gh2018}. Models such as NIH's dbGaP (database of Genotypes and Phenotypes) allow researchers to deposit genomic data with access restrictions applied, making data available for qualified academic research but not for the general public or commercial use.

In Brazil, discussions are ongoing about development of a similar model that balances scientific openness with compliance with the LGPD and Biodiversity Law. An emerging recommendation is that genomic data with identifying potential should be made available through a controlled access model in which: (1) researchers request access to the repository; (2) a data review committee evaluates whether the requester has scientific qualifications and legitimate purposes; (3) the researcher agrees to terms of use restricting use to non-commercial research purposes; and (4) data are provided in a secure computing environment where the researcher can perform analyses but cannot export raw data.

\section{Practical Compliance: A Guide for Researchers}

\subsection{Research Planning with Legal and Ethical Considerations}

Researchers in phylogenomics planning projects involving genomic data should implement procedures to ensure legal and ethical compliance. A recommended practical protocol includes:

\begin{enumerate}
    \item \textbf{Identify data sources:} Document the origin of all genomic data. Are they public data from repositories like GenBank? Is data newly collected? If newly collected, in which country? From which organism?

    \item \textbf{Identify applicable regulatory requirements:} Based on the origin of the data, identify which laws apply. If data are of Brazilian origin, the Biodiversity Law and SisGen apply. If data include humans, Brazilian LGPD or European GDPR apply. If data include indigenous traditional knowledge, UNDRIP and Biodiversity Law apply.

    \item \textbf{Obtain consent and approval:} If research involves new collection of organisms or human data, obtain appropriate consent. Submit protocol for approval from relevant ethics committees (CEUA for animal research, ethics committee for human data).

    \item \textbf{Register with SisGen (Brazil):} If research involves Brazilian biodiversity, register access with SisGen before beginning collection or analysis.

    \item \textbf{Document compliance:} Keep detailed records of all compliance actions --- dates of consent, committee approvals, SisGen registrations.

    \item \textbf{Plan data sharing:} Document how data will be shared. If data can be public, choose an appropriate repository. If data are sensitive, establish controlled access. Document reuse restrictions in metadata.

    \item \textbf{Report compliance in publications:} Publications should include a ``Data Availability'' section documenting: origin of data, where data can be accessed, any access restrictions, compliance with consent requirements and regulations.
\end{enumerate}

\subsection{Documentation of Methods for Transparency}

In addition to legal compliance, methodological transparency practices are essential. Genomic studies should clearly document:

\begin{itemize}
    \item Origin of all data, including information about sequencing (which capture protocol, which coverage depth, which sequencing technology)
    \item Inclusion and exclusion criteria for organisms or samples
    \item Information about data quality (amount of missing data, sequence coverage)
    \item Analysis methodology, including software used and parameters
    \item Computational code used, deposited in a public repository (GitHub, GitLab) for reproducibility
    \item Informed consent and ethics approval if applicable
    \item Author contributions in CASRAI CRediT format
\end{itemize}

\section{Future Perspectives: International Harmonization and Genomic Data as a Common Good}

\subsection{Efforts to Harmonize International Regulations}

Recognizing that divergent regulations between jurisdictions hinder international scientific research, there are ongoing efforts to harmonize legal and ethical requirements around genomic data. GA4GH works with global regulatory authorities to develop shared standards. UNESCO has established recommendations on open data for science. Discussions are underway in international forums about how to balance scientific openness with personal data protection.

A critical emerging question is how to treat genomic data from indigenous communities and specific populations as a ``global common good'' --- a resource that belongs to humanity as a whole for research purposes --- while simultaneously respecting the sovereign rights of nation states and the collective rights of indigenous communities over their genetic resources. This tension remains largely unresolved and should be central in future discussions of international science policy.

\section{Conclusion}

The contemporary researcher in phylogenomics and bioinformatics operates in a complex normative context marked by multiple legal frameworks, often in tension with one another. Brazilian Biodiversity Law and SisGen establish requirements for registration of access to biodiversity. Brazilian LGPD protects personal data, including genomic data from humans. Internationally, the Nagoya Protocol establishes sovereign rights of States over genetic resources. European GDPR imposes even more stringent restrictions on genomic data of Europeans. Indigenous communities possess rights under UNDRIP.

Simultaneously, funding agencies require open data sharing, international repositories such as GenBank facilitate public access to sequences, and scientific norms value reproducibility and transparency that demand broad data and code sharing.

Compliance with this complex normative apparatus requires: (1) knowledge of the multiple applicable legal frameworks; (2) careful planning from the initial project design; (3) detailed documentation of data origin, consent, and compliance; and (4) respectful and participatory engagement with indigenous and traditional communities.

Researchers who successfully navigate this complex normative landscape contribute not only to scientific progress but also to responsible science practice that respects individual rights, community rights, national sovereignty, and the global imperative for scientific sharing for human well-being. The path forward requires continued international collaboration to harmonize regulations, development of responsible access infrastructure, and commitment from the scientific community to ethical practices.

\clearpage
\begin{revise}
\begin{enumerate}
    \item CBD (Rio, 1992, in force 1993): state sovereignty over genetic resources; access requires prior informed consent (PIC); fair sharing of benefits \parencite{cbd1992}.
    \item Nagoya Protocol (2010, in force 2014): operationalizes CBD; requires disclosure of origin in IP; mutually agreed terms; monitoring along the chain \parencite{nagoyaprotocol2014}.
    \item Biodiversity Law (Lei 13.123/2015) + Decree 8.772/2016: mandatory registration in \textbf{SisGen} before accessing Brazilian biodiversity; fine up to R\$200\,thousand; retroactive reclassification if research generates economic product.
    \item LGPD (Lei 13.709/2018, in force 2020): human genomic data = \textbf{sensitive} personal data; explicit consent; right of access, correction, and deletion; mandatory breach reporting.
    \item GDPR (EU, in force 2018): more stringent than LGPD; prohibits transfer outside EU without equivalent safeguards; mandatory DPIA; ``right to be forgotten''.
    \item Arouca Law (Lei 11.794/2008): ethical use of animals; mandatory CEUA in every institution; 3Rs principles (Reduction, Replacement, Refinement); humane euthanasia.
    \item UNDRIP: collective rights of indigenous peoples over traditional knowledge and genetic resources; free prior consent; biopiracy = access without benefit.
    \item INSDC (GenBank/ENA/DDBJ): public access in 24\,h; hold of up to 4 months; controlled access for identifying human data (dbGaP).
    \item Open data mandates: NIH, ERC, EPSRC, FAPESP, CNPq; UNESCO 2021 endorsed by Brazil.
    \item GA4GH: standards for federating controlled access genomic databases; ``responsible access'' model (openness with restrictions on re-identification/commercial use).
    \item Practical compliance (7 steps): identify sources $\to$ regulatory requirements $\to$ consent/approval $\to$ SisGen registration $\to$ document compliance $\to$ plan sharing $\to$ report in publications (Data Availability).
\end{enumerate}
\end{revise}

\clearpage

\printbibliography[heading=subbibliography,title={References}]

